\subsection{State-Space Representation --- The Modern Control Formulation}

The transfer function approach described in the previous section provides a powerful mathematical framework for analyzing single-input, single-output linear time-invariant systems. However, this representation has fundamental limitations that become apparent when we consider more complex systems. A transfer function relates only the input and output of a system, treating the internal workings as a ``black box'' that we cannot examine. This means that transfer functions cannot capture information about initial conditions, cannot directly represent systems with multiple inputs and multiple outputs, and provide no insight into what happens inside the system between the input and output. State-space representation addresses all of these limitations by providing a complete description of system dynamics that tracks all internal variables of interest.

The fundamental concept underlying state-space representation is the \emph{state} of a system. The state can be understood as the minimum amount of information that, combined with knowledge of future inputs, completely determines all future behavior of the system. Consider a simple mechanical system consisting of a mass attached to a spring. If we know the mass's position and velocity at some moment in time, and we know what forces will be applied to the mass in the future, we can predict exactly where the mass will be at any future time. The position and velocity together constitute the state of this system, and no smaller set of information would suffice for this prediction. This is the essence of the state concept: the smallest possible collection of variables that captures everything necessary to forecast the system's evolution.

The state variables are collected into a vector called the \emph{state vector}, typically denoted $\mathbf{x}(t)$. If we have $n$ state variables, then $\mathbf{x}(t)$ is an $n \times 1$ column vector. The inputs to the system are collected into the \emph{input vector} $\mathbf{u}(t)$, and the outputs are collected into the \emph{output vector} $\mathbf{y}(t)$. With these definitions, the dynamics of any linear time-invariant system can be expressed through two fundamental equations. The first equation, called the \emph{state equation}, describes how the state changes with time:

\begin{equation}
\dot{\mathbf{x}}(t) = \mathbf{A}\mathbf{x}(t) + \mathbf{B}\mathbf{u}(t)
\end{equation}

where $\dot{\mathbf{x}}(t)$ denotes the time derivative of the state vector, $\mathbf{A}$ is an $n \times n$ matrix called the \emph{system matrix}, and $\mathbf{B}$ is an $n \times m$ matrix called the \emph{input matrix} (where $m$ is the number of inputs). The second equation, called the \emph{output equation}, describes how the measured outputs relate to the state and inputs:

\begin{equation}
\mathbf{y}(t) = \mathbf{C}\mathbf{x}(t) + \mathbf{D}\mathbf{u}(t)
\end{equation}

where $\mathbf{C}$ is a $p \times n$ matrix called the \emph{output matrix} (where $p$ is the number of outputs) and $\mathbf{D}$ is a $p \times m$ matrix called the \emph{feedthrough matrix}. Together, these four matrices $\{\mathbf{A}, \mathbf{B}, \mathbf{C}, \mathbf{D}\}$ completely characterize the system, just as a transfer function $G(s)$ characterizes a single-input, single-output system. The quadruple $(\mathbf{A}, \mathbf{B}, \mathbf{C}, \mathbf{D})$ is often called a \emph{realization} of the system.

The matrices in state-space representation have direct physical interpretations that make them intuitive once we understand their roles. The system matrix $\mathbf{A}$ governs the natural behavior of the system when no inputs are present; it describes how the state variables influence each other's evolution. The input matrix $\mathbf{B}$ describes how the external inputs affect the rate of change of each state variable. The output matrix $\mathbf{C}$ specifies which state variables (or combinations thereof) are available as measurements, and the feedthrough matrix $\mathbf{D}$ accounts for any input that appears directly in the output without first affecting the state. For many physical systems, $\mathbf{D} = \mathbf{0}$, meaning the input has no direct path to the output; all output signals must first pass through the state dynamics. However, there are important exceptions, such as a simple voltage amplifier where the output voltage is directly proportional to the input voltage.

To make these abstract concepts concrete, let us work through a detailed example of converting a physical mechanical system into state-space form. Consider a mass-spring-damper system consisting of a mass $m$ connected to a spring with stiffness $k$ and a damper with damping coefficient $c$. An external force $f(t)$ is applied to the mass. The equation of motion for this system, derived from Newton's second law, is:

\begin{equation}
m\ddot{x}(t) + c\dot{x}(t) + kx(t) = f(t)
\end{equation}

where $x(t)$ is the displacement of the mass from its equilibrium position. To convert this second-order differential equation into state-space form, we must introduce state variables that represent the fundamental energy-storage elements in the system. Mechanical systems store energy in two forms: kinetic energy associated with motion and potential energy associated with position. Therefore, we choose the position and velocity as our state variables:

\begin{subequations}
\begin{align}
x_1(t) &= x(t) \\
x_2(t) &= \dot{x}(t)
\end{align}
\end{subequations}

With these definitions, we can rewrite the second-order differential equation as two first-order equations. From the definition of $x_2(t)$, we immediately have:

\begin{equation}
\dot{x}_1(t) = x_2(t)
\end{equation}

From the original equation of motion, we solve for the highest derivative $\ddot{x}(t)$:

\begin{equation}
\ddot{x}(t) = -\frac{c}{m}\dot{x}(t) - \frac{k}{m}x(t) + \frac{1}{m}f(t)
\end{equation}

Since $\dot{x}_2(t) = \ddot{x}(t)$, we obtain:

\begin{equation}
\dot{x}_2(t) = -\frac{k}{m}x_1(t) - \frac{c}{m}x_2(t) + \frac{1}{m}f(t)
\end{equation}

Now we can write the state equation in matrix form by collecting $\dot{x}_1$ and $\dot{x}_2$:

\begin{equation}
\begin{bmatrix}
\dot{x}_1 \\
\dot{x}_2
\end{bmatrix}
=
\begin{bmatrix}
0 & 1 \\
-\frac{k}{m} & -\frac{c}{m}
\end{bmatrix}
\begin{bmatrix}
x_1 \\
x_2
\end{bmatrix}
+
\begin{bmatrix}
0 \\
\frac{1}{m}
\end{bmatrix}
f(t)
\end{equation}

This gives us $\mathbf{A} = \begin{bmatrix} 0 & 1 \\ -k/m & -c/m \end{bmatrix}$ and $\mathbf{B} = \begin{bmatrix} 0 \\ 1/m \end{bmatrix}$. For the output equation, suppose we are interested in measuring the position of the mass. Then $y(t) = x(t) = x_1(t)$, which gives $\mathbf{C} = \begin{bmatrix} 1 & 0 \end{bmatrix}$ and $\mathbf{D} = 0$. If instead we measure the velocity, we would have $\mathbf{C} = \begin{bmatrix} 0 & 1 \end{bmatrix}$ and $\mathbf{D} = 0$. If we measure both position and velocity, then $\mathbf{y} = \begin{bmatrix} x_1 \\ x_2 \end{bmatrix}$, so $\mathbf{C} = \mathbf{I}_2$ and $\mathbf{D} = \mathbf{0}$.

Let us verify that this state-space representation is correct by computing the transfer function from the state-space equations and comparing it to what we would obtain using traditional methods. Taking the Laplace transform of the state equation with zero initial conditions:

\begin{equation}
s\mathbf{X}(s) = \mathbf{A}\mathbf{X}(s) + \mathbf{B}U(s)
\end{equation}

Solving for $\mathbf{X}(s)$:

\begin{equation}
(s\mathbf{I} - \mathbf{A})\mathbf{X}(s) = \mathbf{B}U(s) \implies \mathbf{X}(s) = (s\mathbf{I} - \mathbf{A})^{-1}\mathbf{B}U(s)
\end{equation}

The output transform is:

\begin{equation}
\mathbf{Y}(s) = \mathbf{C}\mathbf{X}(s) + \mathbf{D}U(s) = \left[\mathbf{C}(s\mathbf{I} - \mathbf{A})^{-1}\mathbf{B} + \mathbf{D}\right]U(s)
\end{equation}

The quantity in brackets is the transfer function matrix $\mathbf{G}(s)$. For our single-input, single-output system, this becomes a scalar transfer function. Computing $(s\mathbf{I} - \mathbf{A})$ for our mass-spring-damper:

\begin{equation}
s\mathbf{I} - \mathbf{A} =
\begin{bmatrix}
s & -1 \\
\frac{k}{m} & s + \frac{c}{m}
\end{bmatrix}
\end{equation}

Taking the inverse and multiplying by $\mathbf{B}$:

\begin{equation}
(s\mathbf{I} - \mathbf{A})^{-1}\mathbf{B} =
\frac{1}{s^2 + \frac{c}{m}s + \frac{k}{m}}
\begin{bmatrix}
s + \frac{c}{m} & 1 \\
-\frac{k}{m} & s
\end{bmatrix}
\begin{bmatrix}
0 \\ \frac{1}{m}
\end{bmatrix}
=
\frac{1}{m\left(s^2 + \frac{c}{m}s + \frac{k}{m}\right)}
\begin{bmatrix}
1 \\ s
\end{bmatrix}
\end{equation}

Multiplying by $\mathbf{C} = \begin{bmatrix} 1 & 0 \end{bmatrix}$ gives:

\begin{equation}
G(s) = \frac{1}{m\left(s^2 + \frac{c}{m}s + \frac{k}{m}\right)} = \frac{1}{ms^2 + cs + k}
\end{equation}

This matches exactly the transfer function we would obtain by taking the Laplace transform of the original differential equation $m\ddot{x} + c\dot{x} + kx = f$, confirming that our state-space representation is correct.

The choice of state variables is not unique, and different choices can lead to different but equivalent state-space representations. In the mass-spring-damper example, we chose position and velocity as state variables, which is a natural choice because these variables directly correspond to the energy storage in the system. However, we could have chosen different variables. For instance, we might select the position and the acceleration as state variables, though this choice is less intuitive because acceleration is not an energy variable. The important requirement is that the chosen state variables, together with the system inputs, must determine all future behavior. Any set of variables satisfying this requirement is valid, though some choices lead to simpler matrices than others.

Electrical systems provide another important class of examples where state-space representation is natural. Consider a series RLC circuit consisting of a resistor $R$, an inductor $L$, and a capacitor $C$ connected in series with a voltage source $v_s(t)$. The voltage across each component is related to the current by $v_R = Ri$, $v_L = L\frac{di}{dt}$, and $v_C = \frac{1}{C}\int i \, dt$. Kirchhoff's voltage law gives:

\begin{equation}
v_s(t) = Ri(t) + L\frac{di}{dt} + v_C(t)
\end{equation}

We must also account for the relationship between capacitor voltage and current:

\begin{equation}
i(t) = C\frac{dv_C}{dt}
\end{equation}

The energy in this circuit is stored in two forms: magnetic energy in the inductor proportional to $Li^2/2$ and electric energy in the capacitor proportional to $Cv_C^2/2$. Therefore, we choose the inductor current and capacitor voltage as our state variables:

\begin{subequations}
\begin{align}
x_1(t) &= v_C(t) \\
x_2(t) &= i(t)
\end{align}
\end{subequations}

From the capacitor equation, we have $\dot{x}_1 = i/C = x_2/C$. From the Kirchhoff equation, solving for the inductor voltage:

\begin{equation}
L\frac{di}{dt} = v_s - Ri - v_C \implies \dot{x}_2 = \frac{v_s}{L} - \frac{R}{L}x_2 - \frac{1}{L}x_1
\end{equation}

The state equation is therefore:

\begin{equation}
\begin{bmatrix}
\dot{x}_1 \\
\dot{x}_2
\end{bmatrix}
=
\begin{bmatrix}
0 & \frac{1}{C} \\
-\frac{1}{L} & -\frac{R}{L}
\end{bmatrix}
\begin{bmatrix}
x_1 \\
x_2
\end{bmatrix}
+
\begin{bmatrix}
0 \\ \frac{1}{L}
\end{bmatrix}
v_s(t)
\end{equation}

If we measure the capacitor voltage as our output, then $\mathbf{C} = \begin{bmatrix} 1 & 0 \end{bmatrix}$ and $\mathbf{D} = 0$. If we measure the resistor voltage, which is $v_R = Rx_2$, then $\mathbf{C} = \begin{bmatrix} 0 & R \end{bmatrix}$ and $\mathbf{D} = 0$. Computing the transfer function from $v_s$ to $v_C$ yields the familiar second-order form $G(s) = \frac{1}{LCs^2 + RCs + 1}$.

The state-space representation provides several significant advantages over transfer function representation. First and most fundamentally, state-space methods naturally handle multiple-input, multiple-output systems. A transfer function for a MIMO system becomes a matrix of transfer functions, and while this is mathematically tractable, the analysis becomes cumbersome and loses the physical insight that state-space matrices provide. With state-space representation, the matrices $\mathbf{A}$, $\mathbf{B}$, $\mathbf{C}$, and $\mathbf{D}$ capture all input-output relationships in a compact form that is amenable to systematic analysis and design.

Second, state-space representation naturally incorporates initial conditions. The solution to the state equation includes both the zero-input response (due to initial conditions) and the zero-state response (due to inputs):

\begin{equation}
\mathbf{x}(t) = e^{\mathbf{A}t}\mathbf{x}(0) + \int_0^t e^{\mathbf{A}(t-\tau)}\mathbf{B}\mathbf{u}(\tau)\,d\tau
\end{equation}

The matrix exponential $e^{\mathbf{A}t}$, which appears prominently in this solution, provides a complete description of the system's natural behavior. Transfer functions, by contrast, implicitly assume zero initial conditions and cannot directly represent the effect of nonzero initial states. This makes state-space methods essential for problems involving disturbances, startup transients, or any situation where initial conditions matter.

Third, state-space representation provides access to the internal behavior of the system. The state variables represent quantities we can track throughout the system, not just at the input and output ports. This internal visibility is crucial for many control design problems. If we are controlling a satellite, we care not just about its attitude but also about its angular velocity and the states of reaction wheels or thrusters. If we are controlling a chemical reactor, we need to track concentrations throughout the reactor, not just at the inlet and outlet. State-space methods make such multi-variable monitoring and control straightforward.

Fourth, state-space methods are particularly well-suited for modern control design techniques. Optimal control, robust control, adaptive control, and state estimation (including the Kalman filter) are all formulated naturally in state-space terms. The modern control revolution of the 1960s, which introduced these powerful techniques, was enabled by the state-space framework. While transfer function methods remain valuable for analysis and design, especially in the frequency domain, state-space methods provide access to the most powerful tools in the control engineer's repertoire.

To illustrate the conversion from transfer function to state-space representation in the opposite direction, consider a third-order system with transfer function:

\begin{equation}
G(s) = \frac{Y(s)}{U(s)} = \frac{3s + 2}{s^3 + 6s^2 + 11s + 6}
\end{equation}

There are infinitely many state-space realizations for this transfer function, each corresponding to a different choice of state variables. The most common choice is the \emph{controllable canonical form}, which is particularly convenient for control design. For a third-order system with transfer function:

\begin{equation}
G(s) = \frac{b_2s^2 + b_1s + b_0}{s^3 + a_2s^2 + a_1s + a_0}
\end{equation}

the controllable canonical form uses the state variables:

\begin{subequations}
\begin{align}
x_1 &= y \\
x_2 &= \dot{y} + a_2y \\
x_3 &= \ddot{y} + a_2\dot{y} + a_1y
\end{align}
\end{subequations}

With this choice, the state equations become:

\begin{subequations}
\begin{align}
\dot{x}_1 &= x_2 - a_2x_1 \\
\dot{x}_2 &= x_3 - a_1x_1 \\
\dot{x}_3 &= -a_0x_1 + b_0u + b_1\dot{u} + b_2\ddot{u}
\end{align}
\end{subequations}

To avoid derivatives of the input, we augment the state vector to include filters on the input. The standard controllable canonical form for strictly proper transfers (where numerator degree is less than denominator degree) has:

\begin{equation}
\mathbf{A} = \begin{bmatrix}
0 & 1 & 0 \\
0 & 0 & 1 \\
-a_0 & -a_1 & -a_2
\end{bmatrix}, \quad
\mathbf{B} = \begin{bmatrix} 0 \\ 0 \\ 1 \end{bmatrix}, \quad
\mathbf{C} = \begin{bmatrix} b_0 & b_1 & b_2 \end{bmatrix}, \quad
\mathbf{D} = 0
\end{equation}

For our specific example, $a_2 = 6$, $a_1 = 11$, $a_0 = 6$, $b_2 = 0$, $b_1 = 3$, $b_0 = 2$:

\begin{equation}
\mathbf{A} = \begin{bmatrix}
0 & 1 & 0 \\
0 & 0 & 1 \\
-6 & -11 & -6
\end{bmatrix}, \quad
\mathbf{B} = \begin{bmatrix} 0 \\ 0 \\ 1 \end{bmatrix}, \quad
\mathbf{C} = \begin{bmatrix} 2 & 3 & 0 \end{bmatrix}, \quad
\mathbf{D} = 0
\end{equation}

We can verify this realization by computing the transfer function:

\begin{equation}
\mathbf{C}(s\mathbf{I} - \mathbf{A})^{-1}\mathbf{B} =
\begin{bmatrix} 2 & 3 & 0 \end{bmatrix}
\begin{bmatrix}
s & -1 & 0 \\
0 & s & -1 \\
6 & 11 & s+6
\end{bmatrix}^{-1}
\begin{bmatrix} 0 \\ 0 \\ 1 \end{bmatrix}
\end{equation}

Computing the inverse matrix and performing the multiplication yields $\frac{3s+2}{s^3+6s^2+11s+6}$, confirming the realization.

Another common canonical form is the \emph{observable canonical form}, which is the dual of the controllable form. It is obtained by choosing state variables related to the output and its derivatives in a way that makes the $\mathbf{C}$ matrix particularly simple. The observable canonical form has:

\begin{equation}
\mathbf{A} = \begin{bmatrix}
0 & 0 & -a_0 \\
1 & 0 & -a_1 \\
0 & 1 & -a_2
\end{bmatrix}, \quad
\mathbf{B} = \begin{bmatrix} b_0 \\ b_1 \\ b_2 \end{bmatrix}, \quad
\mathbf{C} = \begin{bmatrix} 0 & 0 & 1 \end{bmatrix}, \quad
\mathbf{D} = 0
\end{equation}

For our example:

\begin{equation}
\mathbf{A} = \begin{bmatrix}
0 & 0 & -6 \\
1 & 0 & -11 \\
0 & 1 & -6
\end{bmatrix}, \quad
\mathbf{B} = \begin{bmatrix} 2 \\ 3 \\ 0 \end{bmatrix}, \quad
\mathbf{C} = \begin{bmatrix} 0 & 0 & 1 \end{bmatrix}, \quad
\mathbf{D} = 0
\end{equation}

Both realizations describe the same system; they are related by a similarity transformation. The choice between them depends on the application and the design method being used.

The state-space framework also provides elegant tools for analyzing system properties that are difficult or impossible to assess from transfer functions alone. Controllability, which measures whether we can drive the system from any state to any other state using the available inputs, is determined solely by the pair $(\mathbf{A}, \mathbf{B})$. Observability, which measures whether we can determine the internal state from output measurements, is determined solely by the pair $(\mathbf{A}, \mathbf{C})$. These fundamental properties have no analog in scalar transfer function analysis and are essential for determining whether certain control design objectives are achievable. The controllability matrix $\mathcal{C} = \begin{bmatrix} \mathbf{B} & \mathbf{A}\mathbf{B} & \mathbf{A}^2\mathbf{B} & \cdots & \mathbf{A}^{n-1}\mathbf{B} \end{bmatrix}$ must have full rank for controllability, and the observability matrix $\mathcal{O} = \begin{bmatrix} \mathbf{C} \\ \mathbf{C}\mathbf{A} \\ \mathbf{C}\mathbf{A}^2 \\ \vdots \\ \mathbf{C}\mathbf{A}^{n-1} \end{bmatrix}$ must have full rank for observability.

In summary, state-space representation provides a complete description of linear time-invariant systems that captures all internal dynamics, naturally handles multiple inputs and outputs, incorporates initial conditions explicitly, and enables the full range of modern control design techniques. While transfer functions remain valuable for certain analyses, particularly in the frequency domain, state-space methods are essential for serious control system design and analysis. The matrix equations $\dot{\mathbf{x}} = \mathbf{A}\mathbf{x} + \mathbf{B}\mathbf{u}$ and $\mathbf{y} = \mathbf{C}\mathbf{x} + \mathbf{D}\mathbf{u}$ represent the foundation upon which much of control theory is built.
